\section[Chapter Summary]{Chapter Summary and Transition}
\label{sec:pd_summary}

This chapter established a physics-based framework for the design of a vertical
n-i-p Ge-on-Si waveguide photodetector targeted at O-band reception near
\SI{1310}{\nano\metre}. The discussion progressed from electromagnetic absorption
and the germanium band structure, through drift-diffusion carrier transport and
recombination, to the performance metrics that govern receiver-chain design.
Within this analytical framework, a literature-grounded comparison of detector
architectures identified the waveguide-integrated Ge-on-Si PIN with a U-shaped
anode as the most suitable topology: it preserves process compatibility while
addressing the dominant parasitic-resistance bottleneck that limits conventional
parallel-electrode designs.

The proposed device architecture, specified in \cref{sec:pd_design}, uses an
\SI{8}{\micro\metre}-long by \SI{5}{\micro\metre}-wide Ge absorber with
$W_i=\SI{350}{\nano\metre}$. The U-shaped electrode is represented in the MATLAB
model as a \SI{36}{\percent} reduction of the anode sheet-resistance path
relative to the straight-electrode geometry, while the \SI{103}{\giga\hertz}
response remains the cited Shi~\textit{et al.} benchmark rather than a new
measurement generated by this repository. The coupled FDTD-CHARGE simulation
chain described in \cref{sec:pd_simulation} provides the numerical
infrastructure to evaluate this design at device level and export compact-model
parameters for MATLAB and INTERCONNECT.

\paragraph{Future work: double U-shaped anode and surrogate design-space exploration.}
A natural extension of this work is a validated surrogate model for rapid
design-space exploration. The double U-shaped anode, which would wrap metal
contacts around both long sidewalls of the germanium mesa, is therefore treated
as a design concept and not as a completed simulation result in the current
repository. A full double-U re-optimisation, including the peaking network of
\cref{eq:H_pd_pk}, remains as future work, along with experimental fabrication
and RF characterisation.

The results chapter that follows applies the simulation chain to the O-band
device, presenting the optical generation profiles, absorber length and
intrinsic-thickness sweeps, current-voltage characteristics, bandwidth
estimation, and the system-level sensitivity analysis used to compare the design
against the IEEE~802.3 400GBASE-DR4 receiver target.
